\documentclass[UTF8,11pt,a4paper,fontset=fandol]{ctexart}
\usepackage[margin=23mm]{geometry}
\usepackage{amsmath,amssymb,bm}
\usepackage{newtxtext,newtxmath}
\usepackage{microtype}
\usepackage[hidelinks]{hyperref}
\hypersetup{pdftitle={玻色辅助场行列式权重的可重整性},pdfauthor={}}
\newcommand{\Tr}{\operatorname{Tr}}
\newcommand{\diag}{\operatorname{diag}}
\newcommand{\ee}{\mathrm e}
\newcommand{\ii}{\mathrm i}
\numberwithin{equation}{section}
\setlength{\parskip}{3pt}
\setlength{\emergencystretch}{2em}
\allowdisplaybreaks[1]
\ctexset{section={format=\large\bfseries}}
\pagestyle{empty}
\begin{document}
\begin{center}
{\Large\bfseries 玻色辅助场行列式权重的可重整性}\par
\end{center}
\vspace{3pt}

相互作用可以压制玻色子的高占据态，使物理配分函数有限。
将相互作用作辅助场分解后，某些构型中高占据态的贡献却不再衰减，导致裸迹不收敛。
行列式权重能否在辅助场平均后恢复物理配分函数，取决于截断相关的贡献能否在积分后精确抵消。
下面先给出单位圆边界上的可重整例子及充分判据，再讨论含普通 hopping 的排斥模型。后者的数目守恒传播本征值严格超出单位圆；采用复辅助场路径后，包含高斯因子和 Jacobian 的完整权重处处为正，并准确恢复配分函数。

\section{可重整性的定义}

先讨论 HS 分解没有场依赖标量因子的情形。记单粒子传播矩阵为 $\mathcal B$，相应的多体传播为 $\hat U(\mathcal B)$。
在总粒子数为 $n$ 的子空间中，将该算符记为 $\hat U_n(\mathcal B)$，
它是 $\mathcal B^{\otimes n}$ 在完全对称玻色态上的作用。
若 $\mathcal B$ 可逆，也可写成通常的二次型形式
$\hat U(\mathcal B)=\exp[\hat{\boldsymbol a}^{+}(\log \mathcal B)\hat{\boldsymbol a}]$，
其中 $\hat{\boldsymbol a}$ 收集各模式的湮灭算符。

当固定场迹收敛时，占据数求和给出
\begin{equation}
 \Tr\hat U(\mathcal B)=\prod_j\sum_{n_j=0}^{\infty}r_j^{n_j}
 =\frac1{\det(I-\mathcal B)},
\end{equation}
其中 $r_j$ 为 $\mathcal B$ 的本征值。
当 $|r_j|>1$ 时，占据数项指数增长；当 $|r_j|=1$ 时，占据数项也可能因持续振荡而不收敛。右侧的行列式表达式在这些情况下仍可能有限。
我们称一个辅助场表示可重整，是指该表达式在完整辅助场积分后准确恢复物理配分函数，
\begin{equation}
 Z=\Tr\ee^{-\beta H}
 =\int\mathcal D\phi\,P_{\rm HS}(\phi)\,W(\mathcal B),\qquad W(\mathcal B)=\frac1{\det(I-\mathcal B)}.
 \label{eq:goal}
\end{equation}
这里 $\phi$ 统指完整的辅助场构型，$\mathcal D\phi$ 表示对所有辅助场积分，$P_{\rm HS}(\phi)$ 为归一化的辅助场概率密度。

这一要求也可用截断说明。取与辅助场无关的总粒子数截断 $N\le M$，定义
\begin{equation}
 W_M(\mathcal B)=\sum_{n=0}^{M}\Tr\hat U_n(\mathcal B),\qquad
 Z_M=\Tr_{N\le M}\ee^{-\beta H}.
\end{equation}
固定有限粒子数下的 Hubbard--Stratonovich（HS）恒等式给出 $\int\mathcal D\phi\,P_{\rm HS}(\phi)\,W_M=Z_M$。
若 $Z<\infty$，则可重整性意味着
\begin{equation}
 \lim_{M\to\infty}\int\mathcal D\phi\,P_{\rm HS}(\phi)\,[W_M(\mathcal B)-W(\mathcal B)]=0.
 \label{eq:remainder}
\end{equation}
若 HS 分解产生标量因子，应将它同时保留在截断权重和正规化权重中；改变辅助场积分路径时，还应包括路径 Jacobian。下面先证明一个无需这些附加因子的例子。

\section{排斥 Bose--Hubbard 原子极限中的例子}

考虑单格点上的两个玻色分量 $b,c$，同分量和不同分量之间具有相同的排斥强度 $U>0$。
将化学势项包括在 $H$ 中，
\begin{align}
 H&=\frac U2\bigl[\hat n_b(\hat n_b-1)+\hat n_c(\hat n_c-1)\bigr]
      +U\hat n_b\hat n_c-\mu\hat N \nonumber\\
  &=\frac U2\hat N(\hat N-1)-\mu\hat N,\qquad
 \hat N=\hat n_b+\hat n_c.
 \label{eq:model}
\end{align}
取 $\mu=-3U/2$，则
\begin{equation}
 H=\frac U2\hat N(\hat N+2),\qquad
 Z=\sum_{n=0}^{\infty}(n+1)\ee^{-\beta U n(n+2)/2}<\infty.
 \label{eq:physicalZ}
\end{equation}
这里 $n+1$ 是两个玻色模式在总粒子数为 $n$ 时的态数。总粒子数仍从零取到无穷，没有施加自旋长度约束。

引入通常的 Schwinger 玻色子自旋算符，
\begin{equation}
 \hat S_x=\frac{b^+c+c^+b}{2},\qquad
 \hat S_y=\frac{b^+c-c^+b}{2\ii},\qquad
 \hat S_z=\frac{\hat n_b-\hat n_c}{2}.
\end{equation}
玻色对易关系给出
\begin{equation}
 \hat S_x^2+\hat S_y^2=\hat n_b\hat n_c+\frac{\hat N}{2},
 \qquad
 \hat{\boldsymbol S}^{\,2}=\frac{\hat N(\hat N+2)}4.
\end{equation}
因此该物理哈密顿量恰好是三个自旋平方之和，
\begin{equation}
 H=2U(\hat S_x^2+\hat S_y^2+\hat S_z^2).
 \label{eq:spinsquare}
\end{equation}

现在直接对式~\eqref{eq:spinsquare} 作标准高斯 HS 变换。
令 $\Delta\tau=\beta/L_\tau$，对每个自旋分量使用
\begin{equation}
 \ee^{-2U\Delta\tau\hat S_\alpha^2}
 =\int_{-\infty}^{\infty}\frac{d\phi_\alpha}{\sqrt{2\pi}}\,
    \ee^{-\phi_\alpha^2/2}
    \ee^{\ii\sqrt{4U\Delta\tau}\,\phi_\alpha\hat S_\alpha},
 \qquad \alpha=x,y,z.
 \label{eq:HS}
\end{equation}
将三个分量按 Trotter 分解相乘，并取连续虚时间极限，
\begin{equation}
 \ee^{-\beta H}
 =\lim_{L_\tau\to\infty}
 \left(\ee^{-2U\Delta\tau\hat S_x^2}
       \ee^{-2U\Delta\tau\hat S_y^2}
       \ee^{-2U\Delta\tau\hat S_z^2}\right)^{L_\tau}.
\end{equation}
这一等式在每个固定有限粒子数子空间中成立。
每个时间片只引入式~\eqref{eq:HS} 的三个独立实高斯场。
对应的单粒子传播矩阵为
\begin{equation}
 \mathcal B_{L_\tau}(\phi)
 =\prod_{\ell=1}^{L_\tau}
 \left[
 \ee^{\ii\sqrt{U\Delta\tau}\,\phi_{\ell x}\sigma_x}
 \ee^{\ii\sqrt{U\Delta\tau}\,\phi_{\ell y}\sigma_y}
 \ee^{\ii\sqrt{U\Delta\tau}\,\phi_{\ell z}\sigma_z}
 \right]\in SU(2).
 \label{eq:spinB}
\end{equation}
以下的 $\mathcal B$ 和辅助场积分均指这一标准 HS 表示的连续虚时间极限。

每个构型的本征值都是 $\ee^{\pm\ii\theta}$，可取 $0\le\theta\le\pi$。
在总粒子数为 $n$ 的子空间中，
\begin{equation}
 \Tr\hat U_n(\mathcal B)
 =\sum_{k=0}^n\ee^{\ii(2k-n)\theta}
 =\frac{\sin[(n+1)\theta]}{\sin\theta}.
 \label{eq:character}
\end{equation}
对 $0<\theta<\pi$，这些项随 $n$ 持续振荡，并不趋于零；
在 $\theta=0,\pi$ 处，则分别为 $n+1$ 和 $(-1)^n(n+1)$。
因此每个构型的裸迹都不收敛。本例的发散发生在 $|r_\pm|=1$ 的单位圆边界。
相应的行列式权重为
\begin{equation}
 W(\mathcal B)=\frac1{\det(I-\mathcal B)}
 =\frac1{(1-\ee^{\ii\theta})(1-\ee^{-\ii\theta})}
 =\frac1{4\sin^2(\theta/2)}.
 \label{eq:spinW}
\end{equation}
它非负，在 $\theta=0$ 有极点。下面同时证明这个极点可积，以及完整积分给出正确的 $Z$。

暂时将化学势降低 $\Delta\mu>0$，即在哈密顿量中加入 $\Delta\mu\hat N$，并记
\begin{equation}
 z=\ee^{-\beta\Delta\mu},\qquad 0<z<1.
\end{equation}
这将传播矩阵变成 $z\mathcal B$，其本征值模长均为 $z$。
裸迹此时绝对收敛，故通常的 HS 迹公式给出
\begin{align}
 Z(z)&=\sum_{n=0}^{\infty}(n+1)z^n\ee^{-\beta U n(n+2)/2}
 \nonumber\\
 &=\int\mathcal D\phi\,P_{\rm HS}(\phi)\,
   \frac1{(1-z)^2+4z\sin^2(\theta/2)}.
 \label{eq:spinregulator}
\end{align}
为了撤去化学势偏移，利用逐构型成立的两个不等式，
\begin{equation}
 \frac1{(1-z)^2+4\sin^2(\theta/2)}
 \ \le\
 \frac1{(1-z)^2+4z\sin^2(\theta/2)}
 \ \le\
 \frac1z\,W(\mathcal B).
 \label{eq:spinsandwich}
\end{equation}
让实数 $z$ 从小于一的一侧趋近一。
左侧为非负函数，并逐点单调增加到 $W(\mathcal B)$，所以它的积分也增加到 $\int\mathcal D\phi\,P_{\rm HS}(\phi)\,W(\mathcal B)$。
中间项的积分由式~\eqref{eq:spinregulator} 趋于有限的 $Z$。
第一个不等式因而给出 $\int\mathcal D\phi\,P_{\rm HS}(\phi)\,W(\mathcal B)\le Z$，同时证明了极点可积。
第二个不等式则给出 $zZ(z)\le\int\mathcal D\phi\,P_{\rm HS}(\phi)\,W(\mathcal B)$。
再取同样的极限，得到反向不等式，最终有
\begin{equation}
 \int\mathcal D\phi\,P_{\rm HS}(\phi)\,
 \frac1{\det(I-\mathcal B)}
 =\sum_{n=0}^{\infty}(n+1)\ee^{-\beta U n(n+2)/2}
 =Z.
 \label{eq:exactZ}
\end{equation}

截断项的抵消也随之明确。
对与辅助场无关的总粒子数截断 $N\le M$，有限粒子数 HS 恒等式与式~\eqref{eq:exactZ} 给出
\begin{equation}
 \int\mathcal D\phi\,P_{\rm HS}(\phi)\,[W_M(\mathcal B)-W(\mathcal B)]
 =-\sum_{n=M+1}^{\infty}(n+1)\ee^{-\beta U n(n+2)/2}
 \longrightarrow0.
 \label{eq:tail}
\end{equation}
固定构型中不消失的振荡项，经过相互作用本身的 HS 积分后，成为受排斥能量压制的高占据尾项。
这证明该物理模型的自旋通道 HS 表示满足所定义的可重整性。

这一结论涉及配分函数。若进一步计算密度，应先在 $z<1$ 时求导，再恢复化学势，
\begin{equation}
 \langle\hat N\rangle=\lim_{z\to1}\frac{z}{Z(z)}\frac{dZ(z)}{dz}.
\end{equation}
在单位圆边界，逐构型表达式 $\Tr[\mathcal B(I-\mathcal B)^{-1}]$ 在极点之外恒等于 $-1$。
行列式极点使求导与辅助场积分在这里不能直接交换，因此密度需要保留上述极限次序。

\section{可直接检查的充分判据}

上面的证明可以推广到所有本征值位于单位圆内或边界上的传播。
考虑 $2N_s$ 个模式的粒子数守恒系统，假设物理配分函数有限，且每个固定粒子数子空间的 HS 恒等式成立。
如果每个构型的 $\mathcal B$ 都保持同一个满足 $\Theta^2=-I$ 的 Kramers 时间反演，
并且所有本征值满足 $|r_j|\le1$，则这个 HS 表示可重整。
幺正传播是这一判据的一个特例。

将本征值按 $r_a,r_a^*$ 配对，降低化学势后的权重为
\begin{equation}
 \frac1{\det(I-z\mathcal B)}
 =\prod_{a=1}^{N_s}\frac1{|1-zr_a|^2}\ge0,
 \qquad 0<z<1,\quad |r_a|\le1.
 \label{eq:unitarycriterion}
\end{equation}
固定总粒子数的迹是所有占据分配对应的本征值乘积之和，因此
\begin{equation}
 \left|\Tr\hat U_n(\mathcal B)\right|
 =\left|\sum_{n_1+\cdots+n_{2N_s}=n}\prod_{j=1}^{2N_s}r_j^{n_j}\right|
 \le\binom{n+2N_s-1}{2N_s-1}.
\end{equation}
乘以 $z^n$ 后，占据数求和由 $(1-z)^{-2N_s}$ 控制，可以与辅助场积分交换。
所以 $z<1$ 时仍有 $\int\mathcal D\phi\,P_{\rm HS}(\phi)/\det(I-z\mathcal B)=Z(z)$。

关键的逐构型估计是
\begin{equation}
 |1-zr_a|^2
 =z|1-r_a|^2+(1-z)(1-z|r_a|^2),
\end{equation}
其中 $0\le1-z|r_a|^2\le1$。于是
\begin{equation}
 z|1-r_a|^2\le |1-zr_a|^2\le |1-r_a|^2+1-z,
\end{equation}
从而
\begin{equation}
 \prod_{a=1}^{N_s}\frac1{|1-r_a|^2+1-z}
 \le\frac1{\det(I-z\mathcal B)}
 \le z^{-N_s}\frac1{\det(I-\mathcal B)}.
 \label{eq:spectralbound}
\end{equation}
左侧随 $z$ 从小于一的一侧趋近一而单调增加到目标权重。
积分后取极限，第一个不等式给出行列式积分不超过 $Z$，并保证极点可积；
第二个不等式给出反向不等式。因此行列式积分恰好等于 $Z$。

对于本征值模长可能超过一的传播，还可以利用降低化学势的路径建立另一充分判据。
仍用 $z=\ee^{-\beta\Delta\mu}$，先在 $z$ 足够小时使用收敛迹公式，
再将化学势恢复到原值。

这条恢复化学势的路径可以用单粒子传播矩阵来检查。
考虑任意有限的 $m$ 个模式，假设物理配分函数有限，且每个固定粒子数子空间的 HS 恒等式成立。
若存在与构型无关的常数 $L<\infty$ 和 $\delta>0$，使
\begin{equation}
 \|\mathcal B\|\le L,\qquad
 |\det(I-z\mathcal B)|\ge\delta\quad\text{对所有 }0\le z\le1,
 \label{eq:generalcriterion}
\end{equation}
则该辅助场表示可重整。这里 $\|\mathcal B\|$ 是最大奇异值，也就是矩阵对向量长度的最大放大倍数。上述条件须对概率为一的同一组构型成立。
辅助场按非负、归一化的概率分布积分。

第二个条件检查从 $I$ 到 $I-\mathcal B$ 的整条路径。
沿这条路径，零点只能来自正实本征值 $r=1/z\ge1$。
因此，在传播矩阵一致有界时，一个等价的检查方式是将所有可能的极限构型也包括在内，
确认它们的本征值均避开正实轴区间 $[1,\infty)$。
这一条件允许负实本征值或复本征值的模长超过一。

为证明这一条件足以恢复正确的配分函数，
记固定粒子数的物理配分函数为 $Z_n=\Tr_{N=n}\ee^{-\beta H}$。
比较物理求和与行列式积分，
\begin{equation}
 Z(z)=\sum_{n=0}^{\infty}z^nZ_n,\qquad
 Z_{\rm det}(z)=\int\mathcal D\phi\,P_{\rm HS}(\phi)\,\frac1{\det(I-z\mathcal B)}.
\end{equation}
$m$ 个模式的 $n$ 粒子子空间维数为 $\binom{n+m-1}{m-1}$。
由 $\|\mathcal B\|\le L$，当 $|z|L<1$ 时有
\begin{equation}
 \sum_{n=0}^{\infty}|z|^n\int\mathcal D\phi\,P_{\rm HS}(\phi)\,|\Tr\hat U_n(\mathcal B)|
 \le\sum_{n=0}^{\infty}\binom{n+m-1}{m-1}(|z|L)^n
 =\frac1{(1-|z|L)^m}.
\end{equation}
因此此时可以交换占据数求和与场积分。利用固定有限粒子数下的 HS 恒等式，得到
\begin{equation}
 Z_{\rm det}(z)
 =\int\mathcal D\phi\,P_{\rm HS}(\phi)\sum_{n=0}^{\infty}z^n\Tr\hat U_n(\mathcal B)
 =\sum_{n=0}^{\infty}z^nZ_n=Z(z).
 \label{eq:smallz}
\end{equation}

接下来令实参数 $z$ 从 $0$ 增至 $1$，对应逐渐将化学势升回原值。
由于 $\mathcal B$ 一致有界，$\det(I-z\mathcal B)$ 这个有限阶多项式的系数也一致有界。
结合式~\eqref{eq:generalcriterion} 的非零下界，在实轴线段 $[0,1]$ 周围存在一个对所有构型都适用的复平面区域，满足
\begin{equation}
 |\det(I-z\mathcal B)|\ge\frac\delta2,
 \qquad \left|\frac1{\det(I-z\mathcal B)}\right|\le\frac2\delta.
\end{equation}
这保证 $Z_{\rm det}(z)$ 沿路径保持解析，即在各点附近都能作收敛的 Taylor 展开，且不会因场积分而产生额外奇点。
物理的 $Z(z)$ 在 $|z|<1$ 内也解析，因为 $Z_n\ge0$ 且 $\sum_nZ_n=Z<\infty$。
两个解析函数已由式~\eqref{eq:smallz} 在原点附近相等，
通过相邻的 Taylor 展开可将这一等式连续传递到所有实数 $0\le z<1$。

最后将化学势恢复到原值，即让实数 $z$ 从小于 $1$ 的一侧趋近 $1$。
物理侧的正项级数趋于 $Z$，行列式侧的统一上界允许将极限移入积分。
因此
\begin{equation}
 Z=\lim_{z\to1}Z(z)
 =\lim_{z\to1}Z_{\rm det}(z)
 =\int\mathcal D\phi\,P_{\rm HS}(\phi)\,\frac1{\det(I-\mathcal B)}.
\end{equation}
再结合有限截断的 $\int\mathcal D\phi\,P_{\rm HS}(\phi)\,W_M=Z_M$，便得到式~\eqref{eq:remainder} 中截断余项的积分趋于零。

这个判据把可重整性与恢复化学势过程中行列式的极点位置联系起来。
裸迹的收敛要求控制本征值模长，而上述证明允许本征值越过单位圆，
只要从低化学势回到原模型的路径受到式~\eqref{eq:generalcriterion} 的统一控制。


\section{有 hopping 且传播本征值超出单位圆的模型}
\label{sec:charging_model}

考虑两个格点、每个格点只有一种玻色子的扩展 Bose--Hubbard 模型，
\begin{equation}
 H=K+g\hat N(\hat N+2),\qquad
 K=-t(a_1^+a_2+a_2^+a_1),\qquad
 \hat N=\hat n_1+\hat n_2,\qquad g,t>0.
 \label{eq:charging_dimer}
\end{equation}
将相互作用展开，
\begin{equation}
 H=K+g\sum_{i=1}^2\hat n_i(\hat n_i-1)
       +2g\hat n_1\hat n_2+3g\hat N.
\end{equation}
所以格点内和格点间的密度排斥强度相等，均为 $2g$，化学势为 $\mu=-3g$。
这是含普通 hopping 的物理相互作用模型。
能量下界
\begin{equation}
 H\ge g\hat N^2+(2g-t)\hat N
\end{equation}
保证有限温度的配分函数收敛。

使用无量纲参数 $a=\beta g$、$b=\beta t$。
两个单粒子轨道的能量为 $-t,+t$，因此物理配分函数也可以直接写成
\begin{equation}
 Z=\sum_{n=0}^{\infty}
 \ee^{-a n(n+2)}
 \frac{\sinh[(n+1)b]}{\sinh b}.
 \label{eq:dimer_physical_Z}
\end{equation}
下面从 HS 分解推导同一个配分函数。

将密度平方平移，
\begin{equation}
 H=K+g(\hat N+2)^2-2g\hat N-4g.
\end{equation}
由于 $[K,\hat N]=0$，只需一个实高斯辅助场，且没有 Trotter 误差，
\begin{equation}
 \ee^{-\beta H}
 =\int_{-\infty}^{\infty}\frac{d\phi}{\sqrt{4\pi a}}\,
   \ee^{-\phi^2/(4a)}
   \ee^{4a+2\ii\phi}
   \ee^{(2a+\ii\phi)\hat N-\beta K}.
 \label{eq:dimer_HS}
\end{equation}
它对应的普通数目守恒传播矩阵和标量因子为
\begin{equation}
 \mathcal B(\phi)=\ee^{2a+\ii\phi}\ee^{b\sigma_x},
 \qquad C(\phi)=\ee^{4a+2\ii\phi}=\det\mathcal B(\phi).
 \label{eq:dimer_B}
\end{equation}
两个本征值是 $r_\pm=\ee^{2a\pm b+\ii\phi}$。
当 $2g>t$ 时，它们在每个实场构型中都严格位于单位圆外。

完整正规化权重须保留式~\eqref{eq:dimer_HS} 的标量，
\begin{align}
 W_{\rm reg}(\phi)
 &=\frac{C(\phi)}{\det[I-\mathcal B(\phi)]}
 =\frac1{\det[I-\mathcal B(\phi)^{-1}]}\nonumber\\
 &=\frac1{
 (1-\ee^{-2a+b-\ii\phi})
 (1-\ee^{-2a-b-\ii\phi})}.
 \label{eq:dimer_inverse}
\end{align}
第二步利用了
$\det(I-\mathcal B)=\det\mathcal B\,\det(I-\mathcal B^{-1})$。
逆传播的全部本征值模长小于一，所以右侧可以绝对收敛地展开。

逆传播的多体形式是
$\ee^{-(2a+\ii\phi)\hat N+\beta K}$。
将其占据数求和与高斯积分交换，使用
\begin{equation}
 \int_{-\infty}^{\infty}\frac{d\phi}{\sqrt{4\pi a}}\,
 \ee^{-\phi^2/(4a)}\ee^{-\ii n\phi}=\ee^{-an^2},
\end{equation}
便得到
\begin{equation}
 \int_{-\infty}^{\infty}\frac{d\phi}{\sqrt{4\pi a}}\,
 \ee^{-\phi^2/(4a)}W_{\rm reg}(\phi)
 =\Tr\ee^{-\beta[-K+g\hat N(\hat N+2)]}
 =Z.
 \label{eq:dimer_exact_real}
\end{equation}
最后一步由格点变换 $a_2\mapsto-a_2$ 保证：它将 hopping 变号，同时保持全部密度项。
因此该物理 HS 表示的配分函数准确性已经成立。

实场上的 $W_{\rm reg}$ 一般为复数。
下一节构造复辅助场积分路径，使包含高斯因子、$C(\phi)$ 和路径 Jacobian 的完整采样权重处处为正，同时保留上述配分函数等式。
采用的参数范围为
\begin{equation}
 0<t\le g,\qquad \beta g\ge4.
 \label{eq:contour_parameters}
\end{equation}

\section{完整权重为正的辅助场积分路径}
\label{sec:positive_contour}

$\hat N+2$ 的本征值为整数，因此式~\eqref{eq:dimer_HS} 中的场依赖算符具有 $2\pi$ 周期。
可以把实高斯场准确折叠到 $[0,2\pi)$，其概率密度为
\begin{align}
 P_a(\phi)
 &=\frac1{\sqrt{4\pi a}}\sum_{k\in\mathbb Z}
   \exp\!\left[-\frac{(\phi+2\pi k)^2}{4a}\right]\nonumber\\
 &=\frac1{2\pi}\sum_{n\in\mathbb Z}\ee^{-an^2}\ee^{\ii n\phi}.
 \label{eq:wrapped_gaussian}
\end{align}
实轴上 $P_a>0$、$\int_0^{2\pi}P_a\,d\phi=1$。
它准确保留高斯 HS 恒等式，
\begin{equation}
 \int_0^{2\pi}d\phi\,P_a(\phi)\ee^{\ii n\phi}
 =\ee^{-an^2},\qquad n\in\mathbb Z.
\end{equation}
式~\eqref{eq:wrapped_gaussian} 也给出 $P_a$ 对复数 $\phi$ 的解析延拓。

记完整的被积函数为
\begin{equation}
 f(\phi)=P_a(\phi)W_{\rm reg}(\phi),\qquad
 Z=\int_0^{2\pi}f(\phi)\,d\phi.
 \label{eq:contour_integrand}
\end{equation}
希望将路径写成 $\phi(x)=x+\ii y(x)$，$0\le x\le2\pi$，
使累计积分始终为实数。
为此定义
\begin{equation}
 F(\phi)=\int_0^\phi f(w)\,dw,
 \qquad \operatorname{Im}F(x+\ii y(x))=0.
 \label{eq:contour_level}
\end{equation}
在下文的无奇点条带内，$F$ 与积分路径无关。
沿着式~\eqref{eq:contour_level} 求导，得到可以直接求解的路径方程，
\begin{equation}
 y'(x)=
 -\frac{\operatorname{Im}f(x+\ii y(x))}
        {\operatorname{Re}f(x+\ii y(x))},
 \qquad y(0)=0.
 \label{eq:contour_ode}
\end{equation}
它只使用显式已知的 HS 密度和单粒子行列式。

设路径上 $\operatorname{Re}f>0$，则包含 Jacobian 的完整权重为
\begin{align}
 w(x)
 &=P_a(\phi(x))
   \frac{C(\phi(x))}{\det[I-\mathcal B(\phi(x))]}
   [1+\ii y'(x)]\nonumber\\
 &=f(\phi(x))[1+\ii y'(x)]
 =\frac{|f(\phi(x))|^2}{\operatorname{Re}f(\phi(x))}>0.
 \label{eq:positive_full_weight}
\end{align}
这里的正性属于整个采样权重，包含折叠后的高斯因子、HS 标量和路径 Jacobian。
下面证明这条路径确实存在、闭合，并且不会穿过行列式极点。

在参数范围~\eqref{eq:contour_parameters} 中，考察条带
$|\operatorname{Im}\phi|\le1/2$。
式~\eqref{eq:dimer_inverse} 的极点满足
$\operatorname{Re}\phi=2\pi k$、$\operatorname{Im}\phi=2a\pm b$。
它们与实轴的距离至少为 $2a-b\ge a\ge4$，因此整个条带内没有极点。

高斯 Fourier 系数和逆传播本征值都很小，可以逐项估计
\begin{align}
 |2\pi P_a(\phi)-1|
 &\le2\sum_{n=1}^{\infty}\ee^{-an^2}\cosh(n/2)
 <0.042,\\
 |W_{\rm reg}(\phi)-1|
 &\le(1-\ee^{-a+1/2})^{-2}-1
 <0.065.
 \label{eq:strip_bounds}
\end{align}
这两个界给出
\begin{equation}
 \operatorname{Re}[2\pi f(\phi)]
 >1-0.042-0.065(1+0.042)>0.89.
 \label{eq:positive_real_part}
\end{equation}
在实轴上，还可以使用更强的界
\begin{equation}
 |W_{\rm reg}(x)-1|
 \le(1-\ee^{-a})^{-2}-1<0.039.
\end{equation}
利用 $P_a(x)>0$ 及其归一化，对任意 $0\le x\le2\pi$ 都有
\begin{equation}
 |\operatorname{Im}F(x)|
 =\left|\int_0^xP_a(s)\operatorname{Im}W_{\rm reg}(s)\,ds\right|
 <0.039.
 \label{eq:real_axis_phase}
\end{equation}

另一方面，沿虚方向移动时，
\begin{equation}
 \frac{\partial}{\partial y}\operatorname{Im}F(x+\ii y)
 =\operatorname{Re}f(x+\ii y)>\frac{0.89}{2\pi}.
\end{equation}
从实轴移到 $y=1/2$ 或 $y=-1/2$，虚部的变化幅度至少为
$0.89/(4\pi)>0.070$，超过式~\eqref{eq:real_axis_phase} 中的最大偏差。
因此每个 $x$ 都有唯一的 $|y(x)|<1/2$ 使 $\operatorname{Im}F=0$。
这就证明了式~\eqref{eq:contour_ode} 的解始终留在无极点条带内。

$f$ 具有 $2\pi$ 周期，其一个周期的积分由式~\eqref{eq:dimer_exact_real} 等于实数 $Z$，故
$F(\phi+2\pi)=F(\phi)+Z$。
$\operatorname{Im}F$ 也具有同样的周期，所以 $y(x)$ 周期，且
$y(2\pi)=y(0)=0$。
路径变形不改变积分，最终有
\begin{equation}
 \int_0^{2\pi}w(x)\,dx=Z,\qquad w(x)>0.
 \label{eq:positive_contour_exact}
\end{equation}

再检查原始数目守恒传播。
复路径上的两个本征值满足
\begin{equation}
 |r_\pm(x)|=\ee^{2a\pm b-y(x)}
 \ge\ee^{a-1/2}>1.
 \label{eq:contour_outside}
\end{equation}
因此每个采样构型都具有单位圆外本征值。
固定总粒子数为 $n$ 的裸迹为
\begin{equation}
 \Tr\hat U_n(\mathcal B(\phi(x)))
 =\ee^{[2a-y(x)+\ii x]n}
   \frac{\sinh[(n+1)b]}{\sinh b}.
\end{equation}
其模长随 $n$ 指数增长，裸占据数求和确实发散。

最后，有限总粒子数截断的 HS 被积函数是无奇点的周期解析函数，
所以同一复路径也保持有限截断恒等式。
记
\begin{equation}
 w_M(x)=P_a(\phi(x))\,\phi'(x)\,C(\phi(x))
        \sum_{n=0}^{M}\Tr\hat U_n(\mathcal B(\phi(x))).
\end{equation}
则 $\int_0^{2\pi}w_M(x)\,dx=Z_M$。
结合式~\eqref{eq:positive_contour_exact}，
\begin{equation}
 \lim_{M\to\infty}\int_0^{2\pi}
       [w_M(x)-w(x)]\,dx
 =\lim_{M\to\infty}(Z_M-Z)=0.
 \label{eq:contour_remainder}
\end{equation}
这直接给出截断相关贡献在辅助场积分后精确抵消的物理例子，同时完整采样权重处处为正。

\section{判据及其向更大晶格的推广}
\label{sec:inverse_criterion}

上述推导提供了两个可分开检查的充分条件。
首先，若 $\mathcal B^{-1}$ 的迹一致绝对收敛，可以利用
\begin{equation}
 \frac{C(\phi)}{\det(I-\mathcal B)}
 =\frac{(-1)^m C(\phi)/\det\mathcal B}
        {\det(I-\mathcal B^{-1})}
 \label{eq:inverse_criterion}
\end{equation}
计算正规化积分。
这里 $m$ 为单粒子模式数。
当右侧的 HS 积分恢复一个与原模型等谱的哈密顿量时，配分函数可重整性随即成立。
本例中 $m=2$、$C=\det\mathcal B$，逆传播只把 hopping 变号，而二分晶格变换保证能谱相同。

其次，完整积分的正性可以通过复路径单独建立。
对一个 $2\pi$ 周期、条带内无奇点的被积函数 $f$，若其实轴积分为正实数，且存在 $c,\delta>0$ 满足
\begin{equation}
 \operatorname{Re}f(x+\ii y)\ge c
 \quad (|y|\le\delta),\qquad
 \sup_{0\le x\le2\pi}
 \left|\operatorname{Im}\int_0^x f(s)\,ds\right|<c\delta,
 \label{eq:contour_criterion}
\end{equation}
则上一节的证明直接给出周期路径 $|y(x)|<\delta$，
其完整权重为 $|f|^2/\operatorname{Re}f>0$。
因此可先用逆传播检查配分函数，再用式~\eqref{eq:contour_criterion} 检查正路径的存在性。

这两个条件适用于任意含偶数个格点的有限二分晶格。
每个格点取一种玻色子，考虑
\begin{equation}
 H=K+g\hat N(\hat N+N_s),\qquad
 K=\sum_{ij}a_i^+h_{ij}a_j,\qquad
 \hat N=\sum_i\hat n_i,\qquad N_s\ \text{为偶数}.
 \label{eq:charging_lattice}
\end{equation}
$h$ 为只连接两个子晶格的 Hermitian hopping 矩阵。
相互作用依赖总粒子数，是整体充电能；展开后包含所有格点对之间的密度排斥。
普通最近邻 hopping 可以保留。
物理配分函数有限，因为
\begin{equation}
 H\ge g\hat N^2+(gN_s-\|h\|)\hat N.
\end{equation}

使用恒等式
$g\hat N(\hat N+N_s)=g(\hat N+N_s)^2-gN_s\hat N-gN_s^2$，
同样得到一个高斯辅助场，
\begin{equation}
 \mathcal B(\phi)=\ee^{aN_s+\ii\phi}\ee^{-\beta h},
 \qquad C(\phi)=\ee^{aN_s^2+\ii N_s\phi}
              =\det\mathcal B(\phi),
 \qquad a=\beta g.
\end{equation}
这里使用了 $\Tr h=0$。
由于 $N_s$ 为偶数，
\begin{equation}
 W_{\rm reg}(\phi)
 =\frac1{\det[I-\ee^{-aN_s-\ii\phi}\ee^{\beta h}]}.
 \label{eq:lattice_inverse_weight}
\end{equation}
只要 $gN_s>\|h\|$，逆传播的迹一致绝对收敛。
其 HS 积分恢复 $-K+g\hat N(\hat N+N_s)$。
在一个子晶格上作 $a_i\mapsto-a_i$，便将 $-K$ 变回 $K$，因此正规化积分准确等于 $Z$。

一个直接保证正路径存在的参数范围是
\begin{equation}
 \beta g\ge4,\qquad \|h\|\le\frac{gN_s}{2}.
 \label{eq:lattice_contour_domain}
\end{equation}
设 $h$ 的本征值为 $\epsilon_j$，则逆传播中的正系数
$\ee^{-aN_s+\beta\epsilon_j}$ 满足
\begin{equation}
 \max_j\ee^{-aN_s+\beta\epsilon_j}\le\ee^{-aN_s/2},
 \qquad
 \sum_j\ee^{-aN_s+\beta\epsilon_j}
 \le N_s\ee^{-aN_s/2}\le2\ee^{-a}.
\end{equation}
于是条带内的行列式因子可用几何级数估计，
\begin{equation}
 |W_{\rm reg}(\phi)-1|
 \le
 \exp\!\left[
 \frac{2\ee^{-a+1/2}}{1-\ee^{-a+1/2}}
 \right]-1<0.065,
 \qquad |\operatorname{Im}\phi|\le\frac12.
\end{equation}
在实轴上，将上式中的 $1/2$ 去掉，就得到小于 $0.039$ 的界。
$P_a$ 的界限不变，所以式~\eqref{eq:positive_real_part}--\eqref{eq:real_axis_phase}
以及正路径的证明全部适用。

沿这条路径，原传播的每个本征值都有
\begin{equation}
 |r_j(x)|=\ee^{aN_s-y(x)-\beta\epsilon_j}
 \ge\ee^{aN_s/2-1/2}>1.
\end{equation}
因此式~\eqref{eq:charging_lattice} 给出一个含普通 hopping 的晶格参数族：
数目守恒传播的本征值超出单位圆，完整复路径权重为正，
而截断余项积分仍按式~\eqref{eq:contour_remainder} 趋于零。

\section{怎样判断具体模型}

判别对象是给定哈密顿量及其 HS 表示。
首先检查无限占据空间中的物理配分函数是否有限，再检查辅助场传播矩阵。
本征值满足 $|r_j|\le1$ 且具有 Kramers 配对的表示可以使用式~\eqref{eq:unitarycriterion} 的判据。本征值模长可能超过一时，可以检查式~\eqref{eq:generalcriterion} 的解析性条件，或使用式~\eqref{eq:inverse_criterion} 的逆传播方法；允许复辅助场路径时，再检查式~\eqref{eq:contour_criterion} 的正性条件。
要确认不可重整性，则可以利用传播矩阵的特殊结构，算出行列式积分与物理配分函数之间的差别。

一个特别容易应用的充分条件是，传播矩阵一致有界，且所有本征值的实部都小于一，并留有统一间隔，
\begin{equation}
 \|\mathcal B\|\le L<\infty,\qquad
 \operatorname{Re}r_j\le1-\eta,\qquad 0<\eta\le1.
 \label{eq:realpartcriterion}
\end{equation}
因为对实数 $0\le z\le1$，
\begin{equation}
 |\det(I-z\mathcal B)|
 =\prod_{j=1}^{m}|1-zr_j|
 \ge\prod_{j=1}^{m}(1-z\operatorname{Re}r_j)
 \ge\eta^m,
\end{equation}
所以式~\eqref{eq:generalcriterion} 自动成立。
这一检查只涉及单粒子传播矩阵，允许本征值的模长大于一。

对于前面的 Bose--Hubbard 原子极限，式~\eqref{eq:spinB} 的每个因子都属于 $SU(2)$，
故完整传播幺正，且保持 $\Theta=\ii\sigma_yK$，其中 $K$ 表示复共轭。
物理配分函数的有限性已由式~\eqref{eq:physicalZ} 给出。
这三点直接满足式~\eqref{eq:unitarycriterion} 的判据，所以自旋通道 HS 表示可重整。
此处允许 $\mathcal B$ 接近 $I$，不需要行列式与零之间有统一间隔。

再考虑 $2\times2$ 的 $J_1$--$J_2$ Schwinger 玻色子模型，
\begin{equation}
 H(\mu)=\sum_{\langle ij\rangle}J_{ij}\hat{\boldsymbol S}_i\cdot\hat{\boldsymbol S}_j
       +\lambda\sum_i\hat{\boldsymbol S}_i^2-\mu\sum_i\hat n_i,
 \qquad \hat{\boldsymbol S}_i^2=\frac{\hat n_i(\hat n_i+2)}4.
\end{equation}
这里每个格点有两个玻色分量，$\hat n_i$ 为它们的总占据数，
$J_{ij}$ 取最近邻的 $J_1$ 或次近邻的 $J_2$。
取 $J_1=1$、$0\le J_2\le1$、$\lambda=\mu=16$。
相互作用可写为正系数的自旋平方项，局域系数 $\lambda-2(J_1+J_2)>0$，
从而压制大占据数，物理配分函数有限。
其连续虚时间 HS 表示具有
\begin{equation}
 \mathcal B(\phi;\mu)=\bigoplus_{i=1}^{N_s}\ee^{\beta\mu}U_i(\phi),
 \qquad U_i\in SU(2),\qquad N_s=4.
 \label{eq:rotationmodel}
\end{equation}
旋转矩阵及辅助场概率分布均与 $\mu$ 无关。

对任意 $U_i$，设本征值为 $\ee^{\pm\ii\theta_i}$，则
\begin{align}
 \det(I-\ee^{\beta\mu}U_i)
 &=1-2\ee^{\beta\mu}\cos\theta_i+\ee^{2\beta\mu}\nonumber\\
 &=\ee^{2\beta\mu}\det(I-\ee^{-\beta\mu}U_i).
\end{align}
因此每个构型都满足
\begin{equation}
 W(\phi;\mu)=\ee^{-2N_s\beta\mu}W(\phi;-\mu).
\end{equation}
当 $\mu>0$ 时，右侧的负化学势传播本征值模长为 $\ee^{-\beta\mu}<1$，其迹收敛。
完成辅助场积分，便得到
\begin{equation}
 Z_{\rm det}(\mu)=\ee^{-2N_s\beta\mu}Z(-\mu)<Z(\mu).
 \label{eq:rotationfailure}
\end{equation}
最后的不等式使用了物理配分函数随化学势单调增加，以及 $\ee^{-2N_s\beta\mu}<1$。
这严格证明该表示在正化学势下不可重整。
同一判断适用于所有具有式~\eqref{eq:rotationmodel} 结构、旋转场分布与化学势无关且物理配分函数有限的模型。
对这类表示，负化学势时的裸迹收敛；在 $\mu=0$ 且物理配分函数有限时，幺正 Kramers 判据保证可重整；正化学势时则由式~\eqref{eq:rotationfailure} 留下确定的有限差额。

最后考虑双分量 Bose--Hubbard 模型的密度相互作用，
\begin{equation}
 H-\mu\hat N=H_{\rm hop}
 +U_1\sum_i(\hat n_{b,i}-\hat n_{c,i})^2
 +U_2\sum_i(\hat n_{b,i}+\hat n_{c,i})^2
 -\mu\sum_i(\hat n_{b,i}+\hat n_{c,i}).
 \label{eq:densitybenchmark}
\end{equation}
$H_{\rm hop}$ 为两种分量各自的格点间 hopping。
取一个格点上 $n_b=n_c=n$、其余格点为空的 Fock 态。
该态的 hopping 能量期望为零，因而
\begin{equation}
 \langle H-\mu\hat N\rangle=4U_2n^2-2\mu n.
\end{equation}
当 $U_2<0$ 时，该表达式随 $n$ 增大趋于负无穷。
$U_1$ 项在这个方向上为零，有限化学势也无法压制负的二次项。
因此这一无限占据模型的巨正则配分函数发散，已经无法满足本文以有限 $Z$ 为前提的可重整性定义。

对于 $U_1>0$、$U_2<0$ 的连续高斯解耦，行列式权重的积分也可以直接检查。
两个分量的传播互为复共轭，所以
\begin{equation}
 W(\phi)=\frac1{|\det(I-\mathcal B_b(\phi))|^2}.
\end{equation}
实的总密度场和虚的相对密度场可以分别改变传播的模长与相位。
当它们在空间和虚时间上均匀时，单分量传播写成
\begin{equation}
 \mathcal B_b=\ee^{x+\ii y}\ee^{-\beta(h_0-\mu I)},
\end{equation}
其中 $h_0$ 是单分量 hopping 矩阵，$x,y$ 分别是上述两个高斯场的独立实线性组合。
在 hopping 振幅 $t>0$ 的周期 $3\times3$ 三角格上，$\boldsymbol k=0$ 是能量为 $6t$ 的非简并模式。
选择 $x=\beta(6t-\mu)$、$y=0$，恰有这一个模式的传播本征值等于一。
记 $\delta x=x-\beta(6t-\mu)$，展开该本征值对应的行列式因子，
\begin{equation}
 1-\ee^{\delta x+\ii y}\simeq-(\delta x+\ii y),\qquad
 W(\phi)\ \propto\ \frac1{(\delta x)^2+y^2}.
\end{equation}
其余行列式因子在此处非零。其余辅助场的小变化只会移动这个简单零点，
而高斯概率密度在这些有限场值处严格为正。
因此完整积分在零点附近包含
\begin{equation}
 \int_{(\delta x)^2+y^2<\epsilon^2}
       \frac{d(\delta x)\,dy}{(\delta x)^2+y^2}
 =2\pi\int_0^\epsilon\frac{d\rho}{\rho}=\infty.
\end{equation}
所以这一 $U_2<0$ 表示的物理配分函数和正行列式积分都发散。
有限占据截断得到的是截断模型的配分函数，其极限需要与这里的无限占据问题区分。

两格点模型~\eqref{eq:charging_dimer} 及更大的二分晶格充电模型~\eqref{eq:charging_lattice}，
由逆传播与 hopping 变号证明可重整，再由式~\eqref{eq:contour_criterion} 保证完整复路径权重为正。
排斥 Bose--Hubbard 原子极限则由本征值位于单位圆边界上的 Kramers 判据保证可重整。
旋转型 Schwinger 玻色子模型在正化学势下由式~\eqref{eq:rotationfailure} 确认不可重整。
式~\eqref{eq:densitybenchmark} 的吸引总密度项则使无截断配分函数本身发散。

\end{document}
